\chapter*{Penggunaan AI Generatif}
\addcontentsline{toc}{chapter}{Penggunaan AI Generatif}
\markboth{Penggunaan AI Generatif}{Penggunaan AI Generatif}

Sesuai dengan praktik terbaik saat ini untuk sumber daya pendidikan terbuka,
bab ini mengungkapkan cara kecerdasan buatan (AI) generatif digunakan
dalam penyiapan buku ini.

%-------------------------------------------------------------------
\section*{Kepengarangan}
%-------------------------------------------------------------------

Konten matematis, struktur, contoh, latihan, bukti, dan keputusan pedagogis
dalam buku ini merupakan karya Robert Hildebrand dan para kontributor yang
disebutkan dalam bagian Ucapan Terima Kasih serta Sumber dan Atribusi. AI
digunakan sebagai alat bantu penulisan dan produksi, bukan sebagai penulis.
Penulis telah meninjau, memverifikasi, dan menyunting seluruh materi yang
disusun dengan bantuan AI.

Untuk edisi Bahasa Indonesia ini, bantuan penerjemahan dan produksi diberikan
oleh OpenAI Codex gpt-5.6-sol, Ultra atas permintaan pengguna. Sistem AI tersebut
tidak dicantumkan sebagai pengarang; atribusi Robert Hildebrand dan seluruh
kontributor manusia tetap dipertahankan.

%-------------------------------------------------------------------
\section*{Untuk Apa AI Digunakan}
%-------------------------------------------------------------------

Asisten AI generatif (terutama model bahasa besar) digunakan selama
penyiapan buku ini untuk jenis tugas berikut:

\begin{itemize}
  \item \textbf{Penyuntingan dan penyelarasan naskah} atas uraian untuk
    meningkatkan kejelasan, tata bahasa, dan konsistensi ragam bahasa.
  \item \textbf{Penyusunan draf teks baku}, seperti ringkasan kompatibilitas
    lisensi, pemberitahuan atribusi, metadata aksesibilitas, dan paragraf
    pengantar pembuka bab, yang kemudian direvisi oleh penulis.
  \item \textbf{Pembuatan draf deskripsi teks alternatif} untuk gambar,
    yang ditinjau dan disunting oleh penulis demi ketepatan matematis.
  \item \textbf{Bantuan pemformatan dengan \LaTeX{}}: menghasilkan
    kerangka \texttt{tikz}, memperbaiki definisi makro, menata ulang tabel,
    dan menelusuri galat kompilasi.
  \item \textbf{Penyusunan draf komentar kode dan kode contoh} dalam Python
    dan Excel. Semua contoh numerik dan kode diverifikasi secara independen
    dengan menjalankan kode dan memeriksa hasilnya.
  \item \textbf{Bantuan audit lisensi}: memeriksa silang kompatibilitas
    lisensi sumber, menyusun draf pemberitahuan lisensi, serta mengidentifikasi
    gambar atau bagian teks yang perlu diatribusikan ulang atau ditulis ulang.
  \item \textbf{Koreksi naskah dan pemeriksaan konsistensi} notasi yang
    digunakan dalam berbagai bab.
\end{itemize}

%-------------------------------------------------------------------
\section*{Untuk Apa AI Tidak Digunakan}
%-------------------------------------------------------------------

AI generatif tidak digunakan untuk:

\begin{itemize}
  \item Menghasilkan atau mengada-adakan pernyataan, bukti, teorema,
    atau hasil numerik matematis yang tidak diverifikasi secara independen
    oleh penulis.
  \item Menghasilkan gambar akhir yang digunakan dalam buku terbitan tanpa
    peninjauan. Semua gambar TikZ dibuat secara manual atau dihasilkan oleh AI,
    lalu ditinjau dan diperbaiki.
  \item Menghasilkan solusi latihan tanpa verifikasi.
  \item Mengambil keputusan pedagogis mengenai cakupan, urutan, atau
    penekanan. Keputusan tersebut merupakan keputusan penulis.
\end{itemize}

%-------------------------------------------------------------------
\section*{Contoh Jenis Prompt}
%-------------------------------------------------------------------

Untuk memberi pembaca dan peninjau gambaran konkret tentang cara AI
digunakan, berikut adalah pola prompt representatif yang dipakai selama
penyiapan:

\begin{itemize}
  \item ``Tulis ulang paragraf ini agar lebih jelas tanpa mengubah
    konten matematisnya.''
  \item ``Buat deskripsi teks alternatif satu kalimat untuk gambar TikZ
    berikut yang menyampaikan tata letak visual sekaligus makna
    matematisnya.''
  \item ``Verifikasi bahwa sumber CC~BY~3.0 ini dapat dimasukkan ke dalam
    karya CC~BY-SA~4.0 dan jelaskan mekanisme kompatibilitasnya.''
  \item ``Susun draf contoh Python menggunakan paket \texttt{pyomo} yang
    menyelesaikan program linear berikut; jangan menciptakan opsi solver
    yang keberadaannya tidak Anda yakini.''
  \item ``Identifikasi makro LaTeX yang dirujuk tetapi tidak didefinisikan
    dalam berkas ini dan usulkan perbaikannya.''
\end{itemize}

%-------------------------------------------------------------------
\section*{Verifikasi dan Akuntabilitas}
%-------------------------------------------------------------------

Semua kontribusi berbantuan AI ditinjau oleh penulis. Contoh numerik
diverifikasi dengan menjalankan kode dan, jika sesuai, melalui komputasi
independen menggunakan alat seperti \texttt{scipy.optimize} dan
\texttt{Gurobi}. Kebenaran matematis setiap pernyataan, bukti, dan contoh
dalam buku ini merupakan tanggung jawab penulis.

Pembaca yang menduga adanya galat (baik yang muncul akibat AI maupun bukan)
dianjurkan untuk melaporkannya melalui pelacak isu GitHub proyek atau dengan
mengirim surel ke \texttt{open.optimization@gmail.com}. Ralat akan dicatat
dalam riwayat revisi buku.

%-------------------------------------------------------------------
\section*{Standar Pengungkapan}
%-------------------------------------------------------------------

Pengungkapan ini mengikuti prinsip-prinsip yang ditetapkan oleh program
sumber daya pendidikan terbuka Virginia Tech Publishing dan selaras dengan
norma yang berkembang di komunitas OER: bantuan AI dapat diterima dalam
produksi materi pendidikan terbuka apabila (i) penulis manusia tetap
bertanggung jawab atas konten dan ketepatannya, (ii) AI tidak dicantumkan
sebagai penulis, dan (iii) penggunaan AI diungkapkan secara transparan dalam
karya yang diterbitkan.

\cleardoublepage
